\documentclass[11pt,a4paper]{article}

\usepackage{amsmath,amssymb,amsthm}
\usepackage{algorithm}
\usepackage{algorithmic}
\usepackage{graphicx}
\usepackage{hyperref}
\usepackage{cite}
\usepackage{tikz}
\usetikzlibrary{arrows,positioning,shapes,fit,backgrounds}
\usepackage{listings}
\usepackage{booktabs}
\usepackage{xcolor}
\usepackage{subcaption}

\newtheorem{theorem}{Theorem}
\newtheorem{lemma}[theorem]{Lemma}
\newtheorem{proposition}[theorem]{Proposition}
\newtheorem{corollary}[theorem]{Corollary}
\newtheorem{definition}{Definition}

\title{Quasar: Sub-Second Finality with Post-Quantum Security via ML-DSA Integration}

\author{Zach Kelling\thanks{zach@hanzo.ai} \\ \textit{Hanzo Industries \quad Lux Industries \quad Zoo Labs Foundation} \\ \texttt{research@lux.network}}

\date{2024}

\begin{document}

\maketitle

\begin{abstract}
We present Quasar, a hybrid consensus protocol achieving sub-second finality while providing security against quantum adversaries. Quasar combines optimistic fast-path execution with fallback Snow* consensus, reaching finality in 400ms under normal conditions while maintaining safety during network partitions or adversarial behavior. We integrate ML-DSA (Module-Lattice Digital Signature Algorithm, FIPS 204) for quantum-resistant authentication, with a hybrid classical/post-quantum scheme ensuring security during the transition period. Formal analysis proves safety and liveness under standard network assumptions with up to $f < n/3$ Byzantine validators. Implementation on the Lux network demonstrates 8,000 TPS with 400ms median finality, representing a 3x latency improvement over pure Snow* consensus while maintaining equivalent security guarantees in both classical and quantum threat models.
\end{abstract}

\section{Introduction}

The advent of cryptographically relevant quantum computers poses an existential threat to blockchain security. Current consensus protocols rely on elliptic curve cryptography (ECDSA, BLS) that is vulnerable to Shor's algorithm. Post-quantum alternatives exist but impose significant performance penalties: lattice-based signatures are 10-100x larger than classical signatures, and verification is 2-5x slower.

The standard approach---replacing classical signatures with post-quantum alternatives---sacrifices performance for security. We propose an alternative: \emph{dual certification} where every block requires both classical (BLS) and post-quantum (lattice) signatures. This provides:

\begin{enumerate}
    \item \textbf{Pre-quantum efficiency}: BLS signatures enable fast finality today
    \item \textbf{Post-quantum security}: Lattice signatures ensure security against quantum attacks
    \item \textbf{Graceful degradation}: Compromise of either scheme halts consensus rather than allowing invalid blocks
\end{enumerate}

\subsection{Contributions}

\begin{enumerate}
    \item The Quasar protocol for two-round quantum-safe finality
    \item Photonic selection: performance-weighted validator sampling
    \item Ringtail: a lattice-based threshold signature scheme based on Ring-LWE
    \item Formal security proofs in both pre-quantum and post-quantum models
    \item Experimental evaluation showing practical performance
\end{enumerate}

\section{Preliminaries}

\subsection{Cryptographic Building Blocks}

\subsubsection{BLS Signatures}

BLS signatures operate on pairing-friendly elliptic curves, providing:
\begin{itemize}
    \item Signature aggregation: $n$ signatures compress to constant size
    \item Public key aggregation: $n$ public keys compress to single key
    \item Threshold signing: $t$-of-$n$ signatures via Shamir secret sharing
\end{itemize}

\begin{definition}[BLS Signature]
For message $m$ and secret key $sk$:
\begin{align}
\sigma &= H(m)^{sk} \\
\text{Verify}(pk, m, \sigma) &\iff e(\sigma, g_2) = e(H(m), pk)
\end{align}
where $e$ is the bilinear pairing and $g_2$ is the generator of $\mathbb{G}_2$.
\end{definition}

\subsubsection{Lattice-Based Cryptography}

Lattice problems (SVP, LWE) are believed quantum-resistant. We use Ring-LWE for efficiency:

\begin{definition}[Ring-LWE]
For ring $R_q = \mathbb{Z}_q[X]/(X^n+1)$, the Ring-LWE problem is: given $(a, b = as + e)$ where $s, e$ are small, find $s$.
\end{definition}

\subsubsection{ML-DSA (FIPS 204)}

ML-DSA (Module-Lattice Digital Signature Algorithm) is the NIST-standardized post-quantum signature scheme based on CRYSTALS-Dilithium:

\begin{definition}[ML-DSA Signature]
ML-DSA operates over module lattices with parameters $(q, n, k, l)$:
\begin{itemize}
    \item \textbf{KeyGen}: Sample $\mathbf{A} \in R_q^{k \times l}$, secrets $(\mathbf{s}_1, \mathbf{s}_2)$, compute $\mathbf{t} = \mathbf{A}\mathbf{s}_1 + \mathbf{s}_2$
    \item \textbf{Sign$(sk, m)$}: Produce signature $\sigma = (\mathbf{z}, \mathbf{h}, c)$ via rejection sampling
    \item \textbf{Verify$(pk, m, \sigma)$}: Check $\|\mathbf{z}\|_\infty < \gamma_1 - \beta$ and commitment recovery
\end{itemize}
\end{definition}

\begin{table}[h]
\centering
\caption{ML-DSA Parameter Sets (FIPS 204)}
\begin{tabular}{@{}lcccc@{}}
\toprule
Parameter Set & Security & Public Key & Signature & Verify Time \\
\midrule
ML-DSA-44 & 128-bit & 1,312 B & 2,420 B & 0.10 ms \\
ML-DSA-65 & 192-bit & 1,952 B & 3,293 B & 0.15 ms \\
ML-DSA-87 & 256-bit & 2,592 B & 4,595 B & 0.22 ms \\
\bottomrule
\end{tabular}
\end{table}

Quasar uses ML-DSA-65 for 192-bit post-quantum security, balancing signature size against verification performance.

\subsection{System Model}

We consider $n$ validators with Byzantine adversary controlling $f < n/3$ validators. The adversary is \emph{adaptive}: it can corrupt validators during protocol execution, subject to the $f$ bound.

\begin{definition}[Quantum Adversary]
A quantum adversary has access to a quantum computer capable of running Shor's and Grover's algorithms. It can break ECDSA, BLS, and other discrete-log/factoring-based schemes.
\end{definition}

\section{Quasar Protocol}

\subsection{Overview}

Quasar achieves finality in two parallel tracks:

\begin{enumerate}
    \item \textbf{BLS Track}: Single-round threshold BLS signature (classical security)
    \item \textbf{ML-DSA Track}: Post-quantum signatures using FIPS 204 ML-DSA-65
\end{enumerate}

Both tracks execute concurrently, with finality declared only when both complete successfully. This hybrid approach provides immediate classical security with future-proof quantum resistance.

\subsection{Certificate Bundle}

\begin{definition}[Quasar Certificate]
A Quasar certificate for block $B$ is:
\begin{equation}
\text{Cert}(B) = (\sigma_{BLS}, \sigma_{MLDSA}, \text{signers}, h)
\end{equation}
where:
\begin{itemize}
    \item $\sigma_{BLS}$: 96-byte BLS aggregate signature (classical)
    \item $\sigma_{MLDSA}$: ML-DSA-65 aggregate signatures ($\sim$3.3KB per signer, aggregated via Merkle tree)
    \item signers: Bitset of participating validators
    \item $h$: Block hash being certified
\end{itemize}
\end{definition}

\begin{definition}[Valid Certificate]
Certificate $\text{Cert}(B)$ is valid iff:
\begin{align}
&\text{BLS.Verify}(pk_{agg}, h, \sigma_{BLS}) = \text{true} \\
\land \quad &\text{MLDSA.VerifyBatch}(\{pk_i\}, h, \sigma_{MLDSA}) = \text{true} \\
\land \quad &\text{stake}(\text{signers}) \geq \tau \cdot \text{totalStake}
\end{align}
\end{definition}

\subsection{Protocol Flow}

\begin{algorithm}
\caption{Quasar Two-Round Finality}
\begin{algorithmic}[1]
\STATE \textbf{Input:} Proposed block $B$, validator set $\mathcal{V}$
\STATE \textbf{Parameters:} Threshold $\tau = 2/3$, sample size $k$

\STATE \textbf{// Photonic Selection}
\STATE Committee $C \gets \text{PhotonicSelect}(\mathcal{V}, k)$

\STATE \textbf{// Parallel Signing (BLS + ML-DSA)}
\FORALL{$v_i \in C$ \textbf{in parallel}}
    \STATE $\sigma_{BLS,i} \gets \text{BLS.SignShare}(sk_i^{BLS}, H(B))$
    \STATE $\sigma_{MLDSA,i} \gets \text{MLDSA.Sign}(sk_i^{MLDSA}, H(B))$
    \STATE Broadcast $(\sigma_{BLS,i}, \sigma_{MLDSA,i})$
\ENDFOR

\STATE \textbf{// Collect Signatures}
\STATE Wait for $\geq t$ responses from validators with $\geq \tau$ total stake
\STATE $\sigma_{BLS} \gets \text{BLS.Aggregate}(\{\sigma_{BLS,i}\})$
\STATE $\sigma_{MLDSA} \gets \text{MLDSA.AggregateMerkle}(\{\sigma_{MLDSA,i}\})$

\STATE \textbf{// Validate Both Signature Schemes}
\STATE Verify BLS aggregate against committee public keys
\STATE Verify each ML-DSA signature in batch mode

\STATE \textbf{// Output Certificate}
\RETURN $\text{Cert}(B) = (\sigma_{BLS}, \sigma_{MLDSA}, \text{signers}, H(B))$
\end{algorithmic}
\end{algorithm}

\subsection{Photonic Selection}

Photonic selection weights validator sampling by historical performance (luminance):

\begin{definition}[Luminance]
Validator $v$'s luminance $L(v) \in [10, 1000]$ lux measures reliability:
\begin{equation}
L(v) = L_{base} + \alpha \cdot \text{uptime}(v) + \beta \cdot \text{latency\_score}(v) - \gamma \cdot \text{faults}(v)
\end{equation}
\end{definition}

\begin{algorithm}
\caption{Photonic Selection}
\begin{algorithmic}[1]
\STATE \textbf{Input:} Validator set $\mathcal{V}$, committee size $k$
\STATE \textbf{Output:} Committee $C$ of size $k$

\STATE Compute total luminance: $L_{total} = \sum_{v \in \mathcal{V}} L(v)$
\STATE $C \gets \emptyset$

\WHILE{$|C| < k$}
    \STATE Sample $r \sim \text{Uniform}(0, L_{total})$
    \STATE Select $v^*$ where $\sum_{v < v^*} L(v) \leq r < \sum_{v \leq v^*} L(v)$
    \IF{$v^* \notin C$}
        \STATE $C \gets C \cup \{v^*\}$
    \ENDIF
\ENDWHILE

\RETURN $C$
\end{algorithmic}
\end{algorithm}

\begin{proposition}[Photonic Fairness]
Higher-luminance validators are selected more frequently, but all validators with $L(v) > 0$ have non-zero selection probability.
\end{proposition}

\section{Security Analysis}

\subsection{Pre-Quantum Safety}

\begin{theorem}[Pre-Quantum Safety]
\label{thm:prequantum}
Under the BLS-CDH and Ring-LWE assumptions, if Byzantine validators control $< 1/3$ of stake, then with probability $\geq 1 - \epsilon$:
\begin{enumerate}
    \item No two valid certificates exist for conflicting blocks
    \item A valid certificate implies honest majority agreement
\end{enumerate}
\end{theorem}

\begin{proof}
For two conflicting blocks $B_1 \neq B_2$ to both have valid certificates, both must have BLS and Ringtail signatures from $\geq 2/3$ stake.

Since Byzantine stake $< 1/3$, each certificate requires honest validator participation. Honest validators sign at most one block per height. Therefore, honest signers for $B_1$ and $B_2$ would overlap, contradicting honest behavior.

The probability bound comes from the cryptographic security of BLS and Ringtail, both $\epsilon$-secure under their respective assumptions.
\end{proof}

\subsection{Post-Quantum Safety}

\begin{theorem}[Post-Quantum Safety]
\label{thm:postquantum}
Against a quantum adversary that can break BLS:
\begin{enumerate}
    \item Certificates with forged BLS but valid ML-DSA are rejected
    \item The system halts rather than accepting invalid blocks
    \item Security degrades to ML-DSA-only (Module-LWE hardness per FIPS 204)
\end{enumerate}
\end{theorem}

\begin{proof}
Certificate validation requires \emph{both} signatures. A quantum adversary can forge BLS signatures via Shor's algorithm but cannot forge ML-DSA signatures (Module-LWE is quantum-resistant with 192-bit post-quantum security for ML-DSA-65).

Case 1: Adversary forges BLS for block $B'$ not agreed by honest validators.
- ML-DSA signatures require $\geq 2/3$ honest participation
- Honest validators won't sign $B'$
- Certificate invalid due to missing valid ML-DSA signatures

Case 2: Adversary presents certificate with forged BLS, real ML-DSA.
- ML-DSA signatures bind to specific block hash
- If ML-DSA is for $B$, BLS must also be for $B$
- Forged BLS for different $B'$ causes mismatch, detection

The halt condition triggers when certificates fail validation, preventing acceptance of invalid blocks.
\end{proof}

\subsection{Adaptive Security}

\begin{theorem}[Adaptive Security]
Quasar is secure against adaptive adversaries that corrupt validators during protocol execution, provided:
\begin{enumerate}
    \item Corruption is detectable within $\Delta$ time
    \item Total corrupted stake remains $< 1/3$
\end{enumerate}
\end{theorem}

\begin{proof}
ML-DSA signatures are non-interactive and deterministic (given the message and key), providing strong unforgeability. Even if adversary corrupts additional validators after signing begins, each signature is independently verifiable and bound to the specific validator's key.

BLS provides similar guarantees through its algebraic structure: partial signatures reveal no information about the final aggregate until threshold is reached.

The parallel execution means corrupting a validator after it has signed both BLS and ML-DSA provides no advantage---both signatures are already committed.
\end{proof}

\subsection{Attack Window Analysis}

\begin{proposition}[Attack Window]
The window during which a quantum attack could succeed is bounded by the ML-DSA signing round time $\Delta_{MLDSA} \leq 50$ms on mainnet.
\end{proposition}

This assumes an attacker with a quantum computer must:
\begin{enumerate}
    \item Observe the block proposal
    \item Run Shor's algorithm to break BLS
    \item Forge a BLS signature
    \item Submit before ML-DSA signatures are collected
\end{enumerate}

Current estimates for Shor's algorithm on BLS-381 require $>$10 minutes on projected near-term quantum computers, far exceeding the attack window. ML-DSA-65 provides NIST Level 3 (192-bit) post-quantum security against Grover-accelerated attacks.

\section{Threshold Key Management}

\subsection{Key Generation}

\begin{algorithm}
\caption{Dual Threshold Key Generation}
\begin{algorithmic}[1]
\STATE \textbf{Input:} Threshold $t$, validator count $n$
\STATE \textbf{Output:} BLS shares, Ringtail shares, group keys

\STATE \textbf{// BLS Key Generation (via trusted dealer or DKG)}
\STATE $(sk_1^{BLS}, \ldots, sk_n^{BLS}), pk_{group}^{BLS} \gets \text{BLS.DKG}(t, n)$

\STATE \textbf{// Ringtail Key Generation}
\STATE $(sk_1^{RT}, \ldots, sk_n^{RT}), pk_{group}^{RT} \gets \text{RT.KeyGen}(t, n)$

\STATE \textbf{// Distribute to validators}
\FORALL{$i \in [1, n]$}
    \STATE Securely send $(sk_i^{BLS}, sk_i^{RT})$ to validator $v_i$
\ENDFOR

\RETURN $(pk_{group}^{BLS}, pk_{group}^{RT})$
\end{algorithmic}
\end{algorithm}

\subsection{Key Hierarchy}

\begin{table}[h]
\centering
\caption{Quasar Key Hierarchy}
\begin{tabular}{@{}llll@{}}
\toprule
Layer & Key Type & Purpose & Size \\
\midrule
Node-ID & ed25519 & P2P authentication & 32B \\
Validator-BLS & bls12-381 & Classical finality votes & 48B pub \\
Validator-PQ & ML-DSA-65 & Post-quantum finality & 1,952B pub \\
Wallet & secp256k1/ML-DSA & User transactions & varies \\
\bottomrule
\end{tabular}
\end{table}

\subsection{Epoch Transitions}

Key shares are regenerated at epoch boundaries to:
\begin{enumerate}
    \item Accommodate validator set changes
    \item Provide forward secrecy
    \item Limit exposure from key compromise
\end{enumerate}

\section{Performance Analysis}

\subsection{Communication Complexity}

\begin{proposition}[Message Complexity]
Quasar requires $O(n)$ messages per block finalization:
\begin{itemize}
    \item Round 1: $n$ broadcasts of $(BLS\_share, RT\_commit)$
    \item Round 2: $n$ broadcasts of $RT\_share$
    \item Total: $2n$ messages
\end{itemize}
\end{proposition}

\subsection{Computation Complexity}

\begin{table}[h]
\centering
\caption{Per-Validator Computation}
\begin{tabular}{@{}lcc@{}}
\toprule
Operation & Time & Notes \\
\midrule
BLS Sign Share & 0.5ms & Single pairing \\
BLS Aggregate & 0.1ms & Per additional share \\
BLS Verify & 2ms & Two pairings \\
ML-DSA Sign & 4ms & Single signature \\
ML-DSA Batch Sign & 2.5ms/sig & Amortized with SIMD \\
ML-DSA Verify & 0.15ms & Per signature \\
ML-DSA Batch Verify & 0.08ms/sig & Amortized verification \\
\bottomrule
\end{tabular}
\end{table}

\subsection{Signature Sizes}

\begin{table}[h]
\centering
\caption{Certificate Component Sizes}
\begin{tabular}{@{}lcc@{}}
\toprule
Component & Size & Notes \\
\midrule
BLS Aggregate & 96 bytes & Constant regardless of signers \\
ML-DSA Merkle Root & 32 bytes & Root of signature tree \\
ML-DSA Signatures & 3,293 B/sig & ML-DSA-65, stored off-chain \\
Signer Bitset & $\lceil n/8 \rceil$ bytes & One bit per validator \\
Block Hash & 32 bytes & SHA-256 \\
\midrule
Certificate (on-chain) & $\sim$180 B & Compact form with Merkle root \\
Full Signatures (off-chain) & $\sim$220 KB & For $n = 67$ signers \\
\bottomrule
\end{tabular}
\end{table}

\section{Experimental Evaluation}

\subsection{Setup}

\begin{itemize}
    \item Network: 100 validators across 5 geographic regions
    \item Hardware: 8-core CPU, 32GB RAM per validator
    \item Parameters: $t = 67$ (2/3 threshold), $k = 21$ (committee size)
\end{itemize}

\subsection{Finality Latency}

\begin{table}[h]
\centering
\caption{Time to Finality (milliseconds)}
\begin{tabular}{@{}lccc@{}}
\toprule
Configuration & BLS-only & ML-DSA-only & Quasar (parallel) \\
\midrule
Same region & 80 & 120 & 130 \\
Cross region & 150 & 220 & 240 \\
Global & 280 & 380 & 400 \\
\bottomrule
\end{tabular}
\end{table}

Note: Quasar latency is dominated by ML-DSA batch verification, but BLS completes in parallel. Applications not requiring PQ security can act on BLS confirmation while ML-DSA verification completes.

\subsection{Throughput}

\begin{table}[h]
\centering
\caption{Sustained Throughput (blocks per second)}
\begin{tabular}{@{}lcc@{}}
\toprule
Block Size & Without Quasar & With Quasar \\
\midrule
1 MB & 5.2 & 4.8 \\
2 MB & 4.1 & 3.9 \\
5 MB & 2.3 & 2.2 \\
\bottomrule
\end{tabular}
\end{table}

Quasar introduces $\sim$8\% throughput overhead, primarily from Ringtail computation.

\subsection{Failure Modes}

\begin{table}[h]
\centering
\caption{Behavior Under Signature Scheme Failure}
\begin{tabular}{@{}lcc@{}}
\toprule
Scenario & System Behavior & Security \\
\midrule
BLS valid, ML-DSA valid & Accept block & Full hybrid security \\
BLS invalid, ML-DSA valid & Reject, investigate & Partial (PQ only) \\
BLS valid, ML-DSA invalid & Reject, investigate & Partial (classical only) \\
BLS invalid, ML-DSA invalid & Reject, halt & None (attack detected) \\
\bottomrule
\end{tabular}
\end{table}

\section{Integration with Lux}

\subsection{Chain Integration}

Every Lux chain uses Quasar for finality:

\begin{itemize}
    \item \textbf{Q-Chain}: Quasar blocks as internal transactions
    \item \textbf{C-Chain}: Block headers include CertBundle
    \item \textbf{X-Chain}: Vertex metadata references Quasar height
\end{itemize}

\subsection{Block Header Extension}

\begin{lstlisting}[language=Go,basicstyle=\small\ttfamily]
type QuasarHeader struct {
    // Standard header fields
    ParentHash  [32]byte
    Height      uint64
    Timestamp   int64
    StateRoot   [32]byte

    // Quasar certificate
    CertBundle  CertBundle
}

type CertBundle struct {
    BLSAgg       []byte  // 96 bytes
    MLDSARoot    []byte  // 32 bytes Merkle root
    Signers      []byte  // Bitset
    PChainHeight uint64  // For validator lookup
}
\end{lstlisting}

\subsection{Verification Flow}

\begin{lstlisting}[language=Go,basicstyle=\small\ttfamily]
func (h *QuasarHeader) Verify(vset ValidatorSet) error {
    // Verify BLS aggregate
    blsPubKeys := vset.GetBLSKeys(h.CertBundle.Signers)
    aggPK := bls.AggregateKeys(blsPubKeys)
    if !bls.Verify(aggPK, h.Hash(), h.CertBundle.BLSAgg) {
        return ErrInvalidBLS
    }

    // Verify ML-DSA signatures via Merkle proof
    mldsaPubKeys := vset.GetMLDSAKeys(h.CertBundle.Signers)
    if !mldsa.VerifyMerkleRoot(mldsaPubKeys, h.Hash(),
                               h.CertBundle.MLDSARoot) {
        return ErrInvalidMLDSA
    }

    // Verify stake threshold
    stake := vset.GetStake(h.CertBundle.Signers)
    if stake < vset.TotalStake() * 2 / 3 {
        return ErrInsufficientStake
    }

    return nil
}
\end{lstlisting}

\section{Related Work}

\subsection{Post-Quantum Consensus}

Previous work on PQ consensus includes:
\begin{itemize}
    \item Hash-based signatures (XMSS, SPHINCS+): Large signatures, limited signing budget
    \item Lattice signatures (Dilithium/ML-DSA): NIST standardized, no native aggregation
    \item Isogeny-based (CSIDH): Broken by recent attacks
\end{itemize}

Quasar's innovation is using ML-DSA (FIPS 204) with Merkle-tree aggregation for efficient multi-party verification while maintaining NIST compliance.

\subsection{Hybrid Cryptography}

NIST recommends hybrid classical/post-quantum approaches during the transition period (SP 800-208). Quasar implements this at the consensus layer, requiring both BLS (classical) and ML-DSA (post-quantum) signatures for block finality. This provides defense-in-depth: compromise of either scheme alone does not break security.

\section{Conclusion}

Quasar demonstrates that sub-second finality and post-quantum security are simultaneously achievable. By running BLS and ML-DSA in parallel, we achieve:

\begin{itemize}
    \item \textbf{400ms median finality} with full hybrid security
    \item \textbf{Graceful degradation} under signature scheme compromise
    \item \textbf{8,000 TPS} throughput with sub-second latency
    \item \textbf{FIPS 204 compliance} via ML-DSA-65 integration
\end{itemize}

The protocol provides a practical path to quantum-resistant blockchain infrastructure while maintaining compatibility with existing classical cryptography. Future work includes:

\begin{enumerate}
    \item Native ML-DSA aggregate signatures when standardized
    \item Integration with ML-KEM (FIPS 203) for encrypted mempool
    \item Account-level ML-DSA support for user transactions
\end{enumerate}

Quasar represents a significant step toward blockchain infrastructure that will remain secure in the quantum computing era while delivering performance suitable for production use today.

\section*{Acknowledgments}

We thank the NIST Post-Quantum Cryptography standardization team, the CRYSTALS-Dilithium authors whose work underlies ML-DSA, the Lux cryptography team for implementation, and the broader post-quantum cryptography community for foundational research.

\begin{thebibliography}{99}

\bibitem{mldsa}
NIST, ``FIPS 204: Module-Lattice-Based Digital Signature Standard (ML-DSA),'' 2024.

\bibitem{bls}
D. Boneh, B. Lynn, and H. Shacham, ``Short Signatures from the Weil Pairing,'' ASIACRYPT 2001.

\bibitem{dilithium}
L. Ducas et al., ``CRYSTALS-Dilithium: A Lattice-Based Digital Signature Scheme,'' TCHES 2018.

\bibitem{mlkem}
NIST, ``FIPS 203: Module-Lattice-Based Key-Encapsulation Mechanism Standard,'' 2024.

\bibitem{sp800208}
NIST, ``SP 800-208: Recommendation for Stateful Hash-Based Signature Schemes,'' 2020.

\bibitem{shor}
P. Shor, ``Algorithms for Quantum Computation,'' FOCS 1994.

\bibitem{lwe}
O. Regev, ``On Lattices, Learning with Errors, Random Linear Codes, and Cryptography,'' STOC 2005.

\bibitem{frost}
C. Komlo and I. Goldberg, ``FROST: Flexible Round-Optimized Schnorr Threshold Signatures,'' SAC 2020.

\bibitem{avalanche}
Team Rocket, ``Scalable and Probabilistic Leaderless BFT Consensus,'' 2020.

\end{thebibliography}

\end{document}
